
\chapter{Neoclassical Growth Model}

%\iffalse
%\section{RamseyModel 资料}
%\begin{tabular}{lll}
%\href{C:/Users/wenhu/Documents/keyan/EconomicsResearch/MacroeconomicsLecture/007Ramsey/Ramsey1gl.pdf}{Ramsey1gl}&
%\href{C:/Users/wenhu/Documents/keyan/EconomicsResearch/MacroeconomicsLecture/007Ramsey/Ramseyl001.doc}{Ramseyl001}&
%\href{C:/Users/wenhu/Documents/keyan/EconomicsResearch/MacroeconomicsLecture/007Ramsey/Ramseyluomo1.doc}{Ramseyluomo1}\\
%%\hline
%\href{C:/Users/wenhu/Documents/keyan/EconomicsResearch/MacroeconomicsLecture/007Ramsey/Ramseyluomo2.pdf}{Ramseyluomo2}&
%\href{C:/Users/wenhu/Documents/keyan/EconomicsResearch/MacroeconomicsLecture/007Ramsey/RamseyGrowthModel4.pdf}{RamseyGrowthModel4}&
%\href{C:/Users/wenhu/Documents/keyan/EconomicsResearch/MacroeconomicsLecture/007Ramsey/Ramseyluomo3.pdf}{Ramseyluomo3}\\
%%hline
%\href{https://ocw.mit.edu/courses/economics/14-451-macroeconomic-theory-i-spring-2007/lecture-notes/notes_ch_3.pdf}{MIT Prof. George-Marios Angeletos}&
%\href{https://www.sas.upenn.edu/~jesusfv/lecture10_neoclassical.pdf}{Jesús Fernández-Villaverde}&
%\href{}{}\\
%\href{}{}&
%\href{}{}&
%\href{}{}\\
%%hline
%\href{C:/Users/wenhu/Documents/keyan/EconomicsResearch/MacroeconomicsLecture/007Ramsey/Ramseyextension(glt).pdf}{Ramseyextension(glt)}&
%\href{C:/Users/wenhu/Documents/keyan/EconomicsResearch/MacroeconomicsLecture/007Ramsey/Ramseyextension(hdm).pdf}{Ramseyextension(hdm)}&
%\href{C:/Users/wenhu/Documents/keyan/EconomicsResearch/MacroeconomicsLecture/007Ramsey/Ramseyextension.doc}{Ramseyextension}\\
%%hline
%\href{C:/Users/wenhu/Documents/keyan/EconomicsResearch/MacroeconomicsLecture/007Ramsey/Multi_sector_Growth_Models_Theory_and_Application.pdf}{Multi sector Growth Models Theory and Application}&
%\href{http://www.nippelagerlof.com/teaching/4020/LectureNotes_DynMacro_post.pdf}{Nils-Petter Lagerl ̈of}&
%\href{}{}\\
%%\hline
%\href{}{}&
%\href{}{}&
%\href{}{}\\
%%\hline
%\href{}{}&
%\href{C:/Users/wenhu/Documents/keyan/EconomicsResearch/MacroeconomicsLecture/007Ramsey/}{}&
%\href{C:/Users/wenhu/Documents/keyan/EconomicsResearch/MacroeconomicsLecture/007Ramsey/}{}\\
%%hline
%\href{C:/Users/wenhu/Documents/keyan/EconomicsResearch/MacroeconomicsLecture/007Ramsey/}{}&
%\href{C:/Users/wenhu/Documents/keyan/EconomicsResearch/MacroeconomicsLecture/007Ramsey/}{}&
%\href{C:/Users/wenhu/Documents/keyan/EconomicsResearch/MacroeconomicsLecture/007Ramsey/}{}\\
%%hline
%\href{C:/Users/wenhu/Documents/keyan/EconomicsResearch/MacroeconomicsLecture/007Ramsey/Ramsey1928.pdf}{Ramsey(1928)}&
%\href{C:/Users/wenhu/Documents/keyan/EconomicsResearch/MacroeconomicsLecture/007Ramsey/cass(1965).pdf}{cass(1965)}&
%\href{C:/Users/wenhu/Documents/keyan/EconomicsResearch/MacroeconomicsLecture/007Ramsey/koopmans(1965).pdf}{koopmans(1965)}\\
%%hline
%\href{C:/Users/wenhu/Documents/keyan/EconomicsResearch/MacroeconomicsLecture/007Ramsey/euler_and_capital_transition_equantion_in_Ramsey.pdf}{euler and capital transition equantion in Ramsey}&
%\href{C:/Users/wenhu/Documents/keyan/EconomicsResearch/MacroeconomicsLecture/007Ramsey/Ramsey_(EJ_28).pdf}{Ramsey(EJ-28)}&
%\href{C:/Users/wenhu/Documents/keyan/EconomicsResearch/MacroeconomicsLecture/007Ramsey/Ramsey_Meets_Laibson_in_the_Neoclassical_Growth _Model.pdf}{Ramsey Meets Laibson in the Neoclassical Growth   Model}\\
%%hline
%\href{C:/Users/wenhu/Documents/keyan/EconomicsResearch/MacroeconomicsLecture/007Ramsey/Ramsey1927.pdf}{Ramsey1927}&
%\href{C:/Users/wenhu/Documents/keyan/EconomicsResearch/MacroeconomicsLecture/007Ramsey/The_planar_Ramsey_number_PR.pdf}{The planar Ramsey number }&
%\href{C:/Users/wenhu/Documents/keyan/EconomicsResearch/MacroeconomicsLecture/007Ramsey/The_Strength_of_the_Rainbow_Ramsey_Theorem.pdf}{The Strength of the Rainbow RamseyTheorem}\\
%%\hline
%\href{C:/Users/wenhu/Documents/keyan/EconomicsResearch/MacroeconomicsLecture/007Ramsey/}{}&
%\href{C:/Users/wenhu/Documents/keyan/EconomicsResearch/MacroeconomicsLecture/007Ramsey/}{}&
%\href{C:/Users/wenhu/Documents/keyan/EconomicsResearch/MacroeconomicsLecture/007Ramsey/}{}\\
%%\hline
%\href{C:/Users/wenhu/Documents/keyan/EconomicsResearch/MacroeconomicsLecture/007Ramsey/}{}&
%\href{C:/Users/wenhu/Documents/keyan/EconomicsResearch/MacroeconomicsLecture/007Ramsey/}{}&
%\href{C:/Users/wenhu/Documents/keyan/EconomicsResearch/MacroeconomicsLecture/007Ramsey/}{}\\
%%hline
%\href{C:/Users/wenhu/Documents/keyan/EconomicsResearch/MacroeconomicsLecture/007Ramsey/}{}&
%\href{C:/Users/wenhu/Documents/keyan/EconomicsResearch/MacroeconomicsLecture/007Ramsey/}{}&
%\href{C:/Users/wenhu/Documents/keyan/EconomicsResearch/MacroeconomicsLecture/007Ramsey/}{}\\
%%hline
%\href{}{}&
%\href{}{}&
%\href{}{}\\
%%hline
%\end{tabular}
%\fi


\section{起源与发展}

%\paragraph{Neoclassical Growth Model}
\begin{itemize}
%Neoclassical Growth Model
\item Original contribution of Ramsey (1928). That is why some times it isknown as the Ramsey model.
\item Completed by David Cass (1965) and Tjalling Koopmans (1965).That is why some times it is known as the Cass-Koopmans model.
\item William Brock and Leonard Mirman (1972) introduced uncertainty.
\item Finn Kydland and Edward Prescott (1982) used to create the realbusiness cycle research agenda
\end{itemize}

经济社会由无数的家庭和企业构成。所有的家庭都是相同的，具有相同的偏好，可以选择一个代表性家庭代表所有的家庭。家庭来到人世间就是为了追求幸福，家庭的幸福来自三个方面：物质享受、玩耍和拥有金钱。家庭的幸福来自物质享受(消费)，消费的越多，家庭得到的幸福越多，但是额外消费一单位消费带来消费增量是递减的，即边际效用是递减的。家庭玩耍的越多越幸福，幸福的增量是递减。拥有的货币越多也越幸福，但货币的边际效用是不变的。在Ramsey Model 中，我们只讨论家庭选择怎样的消费路径，或者说按照怎样的消费准则，才能从一生的消费中消费或物质享受获得最大化的幸福。暂时不考虑家
庭从劳动选择和持有货币中获得幸福。假设这是一个没有货币的纯交换经济，家庭的劳动全部投入到生产中没有闲暇，劳动供给是外生给定的。
家庭要获取消费，必须用消费衡量的收入去交换消费。在一个没有货币的纯交换经济中，家庭的收入主要包括劳动收入和资产性收入和企业分红。 

\section{家庭}
   
\begin{itemize}
	\item 家庭目标追求家庭幸福最大化
	\item 资源稀缺：有限的收入
	\item 边际效用递减-消费跨期替代和储蓄的基础
\end{itemize}

\subsection{效用函数(幸福函数)}

经济由大量相同的家庭组成。代表性家庭仅从消费中获得幸福，
家庭消费越多获得的幸福感越强，消费的越少幸福感越弱。所以，家庭的瞬时幸福可以表示为关于消费的瞬时效用函数$ u(c) $，家庭幸福是关于消费的增函数$ u'(c)>0$，家庭幸福随着消费的连续增加，每增加一单位消费带来的幸福增量是递减的，即$ u''(c)<0 $,家庭的幸福函数是严格凹的。家庭的人口出生率为指数$ n $。家庭的劳动供给是无弹性（劳动供给是外生给定的）。家庭是无限存活的，无限存活的家庭目标是追求家庭终生幸福水平的最大化。为简化分析，家庭的初始人口标准化为1，标准化后无限存活的家庭的终生效用水平可以表示为，

\begin{gather}
	\int_{0}^{\infty}u(c(t))e^{nt}e^{-\rho t}dt
	\label{家庭终生效用}
\end{gather}

这里，$ \rho $\footnote{$ \rho $类似利率$ r $,及资本的边际效率$ MEC $，其本质是一种贴现率，把未来效用贴现到当前(0期)时的效用。已知当前效用求未来效用是复利和关系，未来效用转化当前效用是一种贴现关系。祥见《微观经济理论-基本原理与拓展(第9版)》第17章\quad 资本市场p459-485。}为时间偏好率，$ 0<\rho<1 $。

\subsection{家庭的预算约束}

家庭的幸福来自于从产品与服务中的消费获得的，家庭要消费必须要有收入，家庭的收入主要来自劳动收入和储蓄收入(资产性收入)，这里暂时不考虑从企业哪里分得红利。在完全竞争市场，企业的收入所有的收入(产出)在企业实现了利润最大化时，都以要素收入的形式分配给了资本和劳动要素。家庭在$ t $时刻的人口为$ L(t) $，实际工资为$ w_t $，所以在$ t $时刻家庭来自劳动的收入为$ w_t L(t) $。在$ t $时刻，家庭总资产为$ A (t)$，资产的实际收益率为$ r(t) $，资产性收入为$ r(t)A(t) $。家庭在$ t $时刻来自劳动和储蓄两方面总收入是：$ w_tL(t)+r(t)A(t) $。家庭在$ t $期的收入用于消费和储蓄购买金融资产。在两期模型中，消费储蓄的选择可以转化为两期消费的跨期选择，家庭最优的跨期选择是家庭无法通过跨期替代来增加自己的总效用。在无限期，家庭消费的最优选择，任意相邻两期都不能同相互替代增加总效用。

家庭一生的总收入

家庭的目标是在一定的资源约束下，选择消费$ (c(t)) $最大化(\ref{家庭终生效用})式表示的家庭终生效用

\begin{gather}
\max_{c(t)}\int_{0}^{\infty}u(c(t))e^{nt}e^{-\rho t}dt
\label{1.1}
\end{gather}

\iffalse
下面讨论家庭目标实现的资源约束条件。家庭在$ t $时刻的收入来自两方面：一种是劳动收入($ L(t)w_t $),$ L(t) $是劳动数量；另一种是资产性收入($ r(t)A(t) $),$ r(t) $表示资产的报酬率。家庭的收入，用于消费和储蓄（购买金融性资产）。
\fi 
上述关系可以表示如下，
\begin{gather}
	C(t)+\dot{A}(t)=r(t)A(t)+L(t)w(t)
\end{gather}
人均资产，$ a(t) =\frac{A(t)}{L(t)}$
\begin{equation}\label{eqn:3}
\begin{split}
a(t) =\frac{A(t)}{L(t)}\\
\dot{a(t)}=\frac{1}{L(t)^2}(\dot{A(t)}L(t)-A(t)\dot{L(t)})\\
\dot{a(t)}=\frac{\dot{A(t)}}{L(t)}+\frac{A(t)}{L(t)}\frac{\dot{L(t)}}{L(t)}\\
\dot{a(t)}=\frac{r(t)A(t)+L(t)w(t)-C(t)}{L(t)}-na(t)\\
\dot{a(t)}	=r(t)a(t)+w(t)-c(t)-na(t)\\
\dot{a(t)}=(r(t)-n)a(t)+w(t)-c(t)
\end{split}
\end{equation}


\begin{equation}
\begin{split}
a(t) &=\frac{A(t)}{L(t)}\\
\dot{a(t)}&=\frac{1}{L(t)^2}(\dot{A(t)}L(t)-A(t)\dot{L(t)})\\
\dot{a(t)}&=\frac{\dot{A(t)}}{L(t)}+\frac{A(t)}{L(t)}\frac{\dot{L(t)}}{L(t)}\\
\dot{a(t)}&=\frac{r(t)A(t)+L(t)w(t)-C(t)}{L(t)}-na(t)\\
\dot{a(t)}	&=r(t)a(t)+w(t)-c(t)-na(t)\\
\dot{a(t)}&=(r(t)-n)a(t)+w(t)-c(t)
\end{split}
\end{equation}

NO-Ponzi-game condition
\begin{equation}
\lim\limits_{t\to \infty}e^{-(R(t)-nt)}	a(t)
\end{equation}


家庭的终生消费支出$ \int_{0}^{t}e^{\int_{0}^{\infty}-(r(\tau)-n)d \tau } c(t)dt$,家庭终生收入等于初始财富$ a(0) $加上未来收入$ \int_{0}^{\infty}e^{-(R(t)-nt)} w(t) dt$,家庭消费不超过家庭收入
\begin{equation}
	\int_{0}^{t}e^{\int_{0}^{\infty}-(r(\tau)-n)d \tau } c(t)dt\le
	a(0)+\int_{0}^{\infty}e^{-(R(t)-nt)} w(t) dt
\end{equation}
\begin{framed}
	\begin{remark}
		household budget constraints
		\begin{equation}
		\begin{split}
		\dot{a(t)}=(r(t)-n)a(t)+w(t)-c(t)
		\\
		[\dot{a(t)}-(r(t)-n)a(t)]e^{\int_{0}^{t}(r(v)-n)dv}
		=\left(w(t)-c(t)\right)e^{\int_{0}^{t}(r(v)-n)dv}
		\\\dfrac{d}{dt}\left[a(t)e^{\int_{0}^{t}(r(v)-n)dv}\right]=\left(w(t)-c(t)\right)e^{\int_{0}^{t}(r(v)-n)dv}
		\\
		\int_{0}^{\infty}d\left[a(t)e^{\int_{0}^{t}(r(v)-n)dv}\right]=
		\int_{0}^{\infty}w(t)e^{\int_{0}^{t}(r(v)-n)dv}dt-\int_{0}^{\infty}c(t)e^{\int_{0}^{t}(r(v)-n)dv}dt
		\\
		\int_{0}^{\infty}c(t)e^{\int_{0}^{t}(r(v)-n)dv}dt=\lim\limits_{t\to \infty}a(t)e^{\int_{0}^{t}(r(v)-n)dv}+a(0)+\int_{0}^{\infty}w(t)e^{\int_{0}^{t}(r(v)-n)dv}dv
		\end{split}
		\end{equation}
		discounting factor
		\begin{gather}
		R(t)=\int_{0}^{t}r(v)dv
		\end{gather}
		
		\begin{gather}
		\int_{0}^{\infty}c(t)e^{-[R(t)-nt]}=\lim\limits_{t\to \infty}a(t)e^{-[R(t)-nt]}+a(0)+\int_{0}^{\infty}w(t)e^{-[R(t)-nt]}dt	
		\end{gather}
		no-ponzi game conditon
		\begin{gather}
		\lim\limits_{t\to \infty}a(t)e^{-[R(t)-nt]}=0
		\end{gather}
		
		\begin{gather}
		\int_{0}^{\infty}c(t)e^{-[R(t)-nt]}=a(0)+\int_{0}^{\infty}w(t)e^{-[R(t)-nt]}dt	
		\end{gather}
	\end{remark}
\end{framed}


\subsection{First-order condtion}

Household problem,

\begin{equation}
\begin{split}
&\max_{c(t)} U=\int_{0}^{\infty}u(c(t))e^{-(\rho -n)t}dt\\
&s.t. ~	\int_{0}^{\infty}c(t)e^{-[R(t)-nt]}dt\le a(0)+\int_{0}^{\infty}w(t)e^{-[R(t)-nt]}dt
\end{split}
\end{equation}


Lagrange Function

\begin{align}
&L(c(t),\lambda)=\int_{0}^{\infty}u(c(t))e^{-(\rho -n)t}dt+\lambda(t)\left(a(0)+\int_{0}^{\infty}w(t)e^{-[R(t)-nt]}dt-\int_{0}^{\infty}c(t)e^{-[R(t)-nt]}dt\right)
\end{align}

F.O.C

\begin{gather}\label{1.14}
\frac{\partial L(c(t),\lambda)}{\partial c(t)}=0 \Rightarrow  u'(c(t))e^{-(\rho -n)t}=\lambda e^{-[R(t)-nt]}\Rightarrow  u'(c(t))=\lambda e^{-[R(t)-\rho t]}
\end{gather}
differentiate equation(\ref{1.14}) with respect to time t

\textbf{Methods 1}


\begin{equation}
\begin{split}
&u'(c(t))=\lambda e^{-[R(t)-\rho t]}
\\
& \texttt{因为} R(t)=\int_{0}^{t}r(v)dv,\textsf{\texttt{所以}} R'(t)=r(t)，\texttt{上式可变为}
\\
&	u''(c(t))\dot{c(t)}=-\lambda e^{-[R(t)-\rho t]}(r(t)-\rho)
\\
&  \texttt{因为}u'(c(t))=\lambda e^{-[R(t)-\rho t]},\texttt{所以，由上式可得}
\\
&	u''(c(t))\dot{c(t)}=-(r(t)-\rho )u'(c(t))
\\
&		-\frac{u''(c(t))c(t)}{u'(c(t))}\frac{\dot{c(t)}}{c(t)}=r(t)-\rho
\\
&	\rho-\frac{u''(c(t))c(t)}{u'(c(t))}\frac{\dot{c(t)}}{c(t)}=r(t)	
\end{split}
\end{equation}
\textbf{Methods 2}
\begin{align}
%\begin{ split}
& u'(c(t))=\lambda e^{-[R(t)-\rho t]}
\\
& \texttt{因为} R(t)=\int_{0}^{t}r(v)dv,\textsf{\texttt{所以}} R'(t)=r(t)，\texttt{上式可变为}
\\
&\frac{d u'(c(t))}{dt}=-\lambda e^{-[R(t)-\rho t]}(r(t)-\rho)
\\
&	\frac{d u'(c(t))}{dt}=-u'(c(t))(r(t)-\rho)
\\
&\rho-\frac{du'(c(t))/dt}{u'(c(t))}=r(t)
,\texttt{或},
\rho-\frac{u''(c(t))c(t)}{u'(c(t))}\frac{\dot{c(t)}}{c(t)}=r(t)
%\end{split}
\end{align}

Euler equantion 

\begin{gather}
	-\frac{u''(c(t))c(t)}{u'(c(t))}\frac{\dot{c(t)}}{c(t)}=r(t)-\rho
\end{gather}

\begin{gather}
	\texttt{令}\quad\theta=	-\frac{u''(c(t))c(t)}{u'(c(t))},\quad \texttt{则}，
\end{gather}

\begin{gather}
	\frac{\dot{c(t)}}{c(t)}=\frac{1}{\theta}(r(t)-\rho)
\end{gather}


\textbf{Methods 3}


\begin{equation}
	\begin{split}
	&\ln u'(c(t))=\ln(\lambda e^{-[R(t)-\rho t]})
	\\
	&\ln u'(c(t))=\ln\lambda-[R(t)-\rho t]
	\\
	&\ln u'(c(t))=\ln\lambda-[R(t)-\rho t]
	\\
	&\texttt{两边对时间t求导}
	\\
	&-\frac{1}{u'(c(t))}u''(c(t))\dot{c(t)}=r(t)-\rho
	\\
	&\rho-\frac{u''(c(t)) c(t)}{u'(c(t))}\frac{\dot{c(t)}}{c(t)}=r(t)
	\end{split}
\end{equation}

\textbf{最优控制理论方法(optimal controll)}

Hanmiltong fuction

\begin{gather}
H=u(c(t))e^{-(\rho-n)t}+\mu(t)[(r(t)-n)a(t)+w(t)-c(t)]
\end{gather}
\footnote{巴罗，萨拉伊玛丁·《经济增长》(第二版)P89.}
F.O.C

\begin{align}\label{1.25}
&\frac{\partial H}{\partial c}=0 \Rightarrow u'(c(t))e^{-(\rho-n)t}=\mu(t)
\end{align}

\begin{align}\label{1.26}
	\frac{\partial H}{\partial a(t)}=-\dot{\mu(t)}
\Rightarrow \mu(t)(r(t)-n)	=-\dot{\mu(t)}
%\Rightarrow \frac{\dot{\mu(t)}}{\mu(t)}=r(t)-n
\end{align}
\begin{gather}
\texttt{对}u'(c(t))e^{-(\rho-n)t}=\mu(t)\texttt{两边对时间t求导}
\\
\dot{u(t)}=\dfrac{d u'(c(t))}{dt}e^{-(\rho-n)t}-(\rho-n)u'(c(t))e^{-(\rho-n)t}	
\end{gather}

把$ \dot{u(t)}=\dfrac{d u'(c(t))}{dt}e^{-(\rho-n)t}-(\rho-n)u'(c(t))e^{-(\rho-n)t} $	和$ u'(c(t))e^{-(\rho-n)t}=\mu(t) $代入$ \mu(t)(r(t)-n)	=-\dot{\mu(t)} $可得，


\begin{gather}
	\mu(t)(r(t)-n)	=-\dot{\mu(t)}
	\\
  u'(c(t))e^{-(\rho-n)t}\left(r(t)-n\right)=\dfrac{d u'(c(t))}{dt}e^{-(\rho-n)t}-(\rho-n)u'(c(t))e^{-(\rho-n)t}
  \\
  u'(c(t))(r(t)-n)=-\dfrac{d u'(c(t))}{dt}+(\rho-n)u'(c(t))
  \\
  u'(c(t))[(r(t)-n)-(\rho-n)]=-\dfrac{d u'(c(t))}{dt}
  \\
  -\dfrac{d u'(c(t))/dt}{u'(c(t))}=r(t)-\rho
  \\
  r=\rho-\left[\dfrac{u''(c(t))\cdot c(t)}{u'(c(t))}\cdot\dfrac{\dot{c(t)}}{c(t)}\right]
\end{gather}

欧拉方程

\begin{gather}\label{欧拉方程(巴罗)}
	r=\rho-\left[\dfrac{u''(c(t))\cdot c(t)}{u'(c(t))}\cdot\dfrac{\dot{c(t)}}{c(t)}\right]
\end{gather}
上式(\ref{欧拉方程(巴罗)})表明家庭选择消费使收益率等于时间偏好率$ \rho $由于人均消费连续增加导致的消费边际效用递减率。
   (\ref{欧拉方程(巴罗)})式左边的利率$ r $是储蓄的收益率。该式的右边可视为消费的收益率。
\begin{align}\label{1.27}
  	TVC \quad \lim\limits_{t\to \infty}\mu(t)a(t)e^{-[R(t)-n]}\ge 0
\end{align}

将$ r=\rho-\left[\dfrac{u''(c(t))\cdot c(t)}{u'(c(t))}\cdot\dfrac{\dot{c(t)}}{c(t)}\right] $变形为，
\begin{gather}
\dfrac{\dot{c(t)}}{c(t)}=-\dfrac{u'(c(t))}{u''(c(t))\cdot c(t)}(r-\rho)	
\end{gather}



\chapter{RCK Model}
\footnote{罗默《高级宏观经济学》第二章}


A部分: Ramsey-Cass-Koopmans模型

\section{ 假设}

\subsection{厂商}
\begin{itemize}
\item 生产函数: Y(t)=F(K(t), A(t)L(t)). 关于F的假设与第一章相同。
\item 要素市场和产出市场都是竞争性的。
\item 厂商利润最大化。家庭拥有厂商（企业）。	
\end{itemize}
 

\subsection{家庭} 

家庭效用最大化
\begin{gather}
	\max\int_{0}^{\infty}e^{-\rho t}u\left(C(t)\right)\dfrac{L(t)}{H}dt
\end{gather}
此处，\\
$ C(t): $ 家庭中每个成员的消费,\\
$ u(.): $ 瞬时效用函数,\\
$ L(t): $ 经济的总人口,$  dL(t)/dt = nL(t) $\\
$ H: $ 家庭的数量,\\
$ u(C(t))L(t)/H: $ 家庭的$t $时刻的 总瞬时效用,\\
$ \rho:  $贴现率\\
$ K(0): $ 初始资本\\


\noindent 瞬时效用函数为：
\begin{gather}
	u\left(C(t)\right)
	=\dfrac{C(t)^{1-\theta}}{1-\theta},\theta>0,
	\rho-n-(1-\theta)g>0.
\end{gather}


\begin{itemize}
\item 常数相对风险厌恶效用函数（Constant-relative-risk-aversion (CRRA) utility function）：相对风险厌恶系数为-$ Cu''(C)/u'(C)=\theta $
\item RRA系数: 0
\item 跨期替代弹性（Elasticity of substitution）:$ 1/\theta $
	
\end{itemize}




补充材料：\\
严格定义：\\
\begin{itemize}
	\item Constantinides（1990）和Weil（1989）按照投资者的值函数来定义投资者的RRA系数，$ RRA\equiv -\dfrac{WJ_{ww}}{J_{W}} $，将RRA系数进一步表示为投资者的值函数关于财富的弹性，得到：
	$ RRA=-\dfrac{dJ_w/J_w}{dW/W} $。因此，投资者的RRA系数的含义是投资者的财富变动1个百分点，投资者的边际效用变动多少个百分点。RRA系数刻画了投资者关于风险的态度。
	\item 投资者的IES系数则定义为当股票的溢价$ \mu-r $保持不变时，经济中利率水平增加1个单位，投资者的最优消费增长率增加的幅度，即如下偏微分：
	$ IES \equiv \dfrac{\partial [E(dC/C)/dt]}{\partial r}|_{\mu-r}$
	\item 当投资者的效用函数为CRRA时，投资者的RRA系数和IES系数互为倒数：$ RRA\times IES＝1 $。当投资者的效用函数不再是CRRA时，以上关系式就不再成立。
\end{itemize}


\textbf{相对风险规避系数的含义：}\\
假设某个投资者获得一笔奖金，并且必须在两种支付方法之间选择一种支付方法。

第一种支付方法是确定的：投资者得到数量确定的$ C-\rho $。

第二种支付方法是不确定的：投资者以50％的概率得到较多的$ C+y $，以50％的概率得到较少的$ C-y $。

那么，多大的$ \rho $可以使得投资者对于这两种支付方式是无差异的？

投资者对于两种支付方式是无差异的，因而$ \rho $的选择使得如下方程成立，
\[u[C-\rho(C,y)] =0.5u(C+y)+0.5u(C-y)\]

其中$ \rho(C,y) $表示$ \rho $是$ C $和$ y $的函数，是严格增加的、严格凹的效用函数。

将$ u[C-\rho(C,y)] $在C点处泰勒展开，得到

将在C点处泰勒展开，得到
$ u[C-\rho(C,y)] \sim u(C)-\rho u'(C)$
将$ u(C+y) $在C点处泰勒展开，得到
\[u(C-y)\simeq u u(C)-yu'(C) +0.5y^2 u''(C)\]

合并，得到
\[\rho(C,y) \sim \dfrac{1}{2}\dfrac{y^2}{C}\left[
-\dfrac{Cu''(C)}{u'(C)}\right]\]


因此，只要投资者是风险规避型的投资者，即$ u''<0 $，那么$ \rho(C,y) $就会大于0。也就是说，不确定性将会降低投资者的效用水平，风险规避型的投资者需要一定数量的补偿。

下面给出CRRA效用函数中的RRA系数的含义。将CRRA型效用函数代入方程（8.5）右边，得到
\[\dfrac{\rho}{y}\sim \dfrac{\alpha}{2}\dfrac{y}{C}\]
对于给定的$ C $和$ y $，$ \rho $越大，$ \alpha $就越大。而更大的$ \rho $意味着投资者在面临同样的不确定性时，宁愿选择更小的确定性支付数量。参数$ \alpha $越大，投资者就惧怕风险。因此，参数$ \alpha $确实刻画了投资者惧怕风险的程度。

通过设定假定的情景（给定$ C $和$ y $），可以调查投资者的$ \rho $是多少。然后就可以估计投资者的相对风险规避系数。大量的统计调查表明一般的美国投资者的$ \alpha $是小于2的。

\section{家庭与厂商的行为}
\subsection{ 厂商}
假设资本没有折旧。厂商最优使用资本和劳动以最大化利润
\begin{gather}
	\pi(K,L)=F(K,AL)-rK-wAL
\end{gather}

$ w  $是单位有效劳动的实际工资。
此规划问题的关于资本的一阶条件为
\begin{gather}
	F_K = r =f'(k)
\end{gather}

关于有效劳动$ AL $的一阶条件为
关于的FOC:
\begin{gather}
F_{AL} = w =f(k) – kf'(k)
\end{gather}

Hint: $ Y=F=ALf(K/AL) $
因此，人均工资为$ W(t)=A(t)w(t) $

\subsection{家庭的预算约束}

家庭将$ r  $和$ w  $视为给定。\\

定义$ R(t)=\int_{0}^{t}r(s)ds $。

\textbf{问题：}0时刻投资一单位产出，在t时获得的回报为$ e^{R(t)} $ 。[why?]\\
\textbf{解答：}将[0,t]分为n个∆t。0时投资1单位产出t时回报为
\[X=(1+r_1\Delta t)\cdots(1+r_n\Delta t)\]
取对数，$ \ln X=∑ln(1+ri∆t)=∑ri∆t=\int_{0}^{t} r(\tau)d\tau$。 因此，$ X=e^{\int_{0}^{t} r(\tau)d\tau} $

家庭的预算约束方程为

\begin{gather}
	\int_{0}^{\infty}e^{-R(t)}C(t)\dfrac{L(t)}{H}dt
	\le \dfrac{K(0)}{H}+\int_{t=0}^{\infty}e^{-R(t)}W(t)
	\dfrac{L(t)}{H}dt
\end{gather}
变形得到
\begin{gather}
	\dfrac{K(0)}{H}+\int_{0}^{\infty}e^{-R(t)}[W(t)-C(t)]\dfrac{L(t)}{H}dt\ge 0
\end{gather}
可以记为

\begin{gather}
	\lim\limits_{s\to \infty}
	\left\{
	\dfrac{K(0)}{H}+\int_{t=0}^{s}e^{-R(t)}
	[
	W(t)-C(t)
	]\dfrac{L(t)}{H}
	\right\}\ge 0
\end{gather}
家庭在s时刻的资本持有数量为
\begin{gather}
	\dfrac{K(s)}{H}=e^{R(s)}\dfrac{K(0)}{H}+\int_{t=0}^{s}e^{R(s)-R(t)}[
	W(t)-C(t)
	]\dfrac{L(t)}{H}
\end{gather}
因此，家庭的预算约束方程可以记为
\begin{gather}
	\lim\limits_{s\to\infty}
	\left\{
	e^{-R(s)}\dfrac{K(s)}{H}\ge 0
	\right\}
\end{gather}
含义：极限中家庭所持有资产的现值不能为负。此为非No-Ponzi-game Condition。

\subsection{家庭的最大化问题} 

定义$  c(t)=C(t)/A(t) $，那么家庭的效用函数可以记为
\begin{align}
	\dfrac{C(t)^{1-\theta}}{1-\theta}&=\dfrac{[A(t)c(t)]^{1-\theta}}{1-\theta}
	\\
	&\dfrac{[A(0)e^{gt}]^{1-\theta}c(t)^{1-\theta}}{1-\theta}
	\\
	&=A(0)^{1-\theta}e^{(1-\theta)gt}\dfrac{c(t)^{1-\theta}}{1-\theta}
\end{align}

由于$ L(t) = L(0)e^{nt} $，所以
\begin{align}
	U&=\int_{t=0}^{\infty}e^{-\rho t}\dfrac{C(t)^{1-\theta}}{1-\theta}
	\dfrac{L(t)}{H}dt
	\\
	&=\int_{t=0}^{\infty}e^{-\rho t}
	\left[
	A(0)^{1-\theta}e^{(1-\theta)gt}\dfrac{c(t)^{1-\theta}}{1-\theta}
	\right]\dfrac{L(0)e^{nt}}{H}dt
	\\
	&=A(0)^{1-\theta}\dfrac{L(0)}{H}\int_{t=0}^{\infty}e^{-\rho t}
	e^{(1-\theta)nt}\dfrac{c(t)^{1-\theta}}{1-\theta}dt
	\\
	&=B\int_{t=0}^{\infty}e^{-\beta t}\dfrac{c(t)^{1-\theta}}{1-\theta}
\end{align}

此处，

\begin{gather}
	B=A(0)^{1-\theta}L(0)/H ,and \beta=\rho-n-(1-\theta)g
\end{gather}
下面把预算约束方程也改写为单位有效劳动的形式。

\begin{gather}
\int_{t=0}^{\infty}e^{-R(t)}C(t)\dfrac{L(t)}{H}dt\le \dfrac{K(0)}{H}+\int_{t=0}^{\infty}e^{-R(t)}W(t)
\dfrac{L(t)}{H}
\\
\Longleftrightarrow
\\
\int_{t=0}^{\infty}e^{-R(t)}c(t)\dfrac{A(t)L(t)}{H}dt
\le
k(0)\dfrac{A(0)L(0)}{H}+\int_{t=0}^{\infty}e^{-R(t)}w(t)\dfrac{A(t)L(t)}{H}dt
\\
\Longleftrightarrow
\\
\int_{t=0}^{\infty}e^{-R(t)}c(t)e^{(n+g)t}dt
\le
k(0)+\int_{t=0}^{\infty}e^{-R(t)}w(t)e^{(n+g)t}dt
\end{gather}



此处使用了$ A(t)L(t) equals A(0)L(0)e^{(n+g)t}$
可以将预算约束的No-Ponzi-game Condition改写为：


\begin{gather}
\lim\limits_{s\to \infty}e^{-R(s)}e^{(n+g)s}k(s)\ge 0
\end{gather}



\subsection{ 家庭行为（家庭的规划问题以及离散化求解方法）}
家庭的效用最大化问题可以表示为如下规划问题：
\begin{gather}
	\max\left\{B\int_{t=0}^{\infty}e^{-\beta t}\dfrac{c(t)^{1-\theta}}{1-\theta}dt
	\right\}
\end{gather}
s.t.

\begin{gather}
	\int_{t=0}^{\infty}e^{-R(t)}c(t)e^{(n+g)t}dt\le
	k(0)+\int_{t=0}^{\infty}e^{-R(t)}w(t)e^{(n+g)t}dt
\end{gather}

如何求解?


此规划问题的目标函数和预算约束方程都是连续时间表达形式。我们采用离散化目标函数和预算约束方程的求解方法。因此，得到如下新的规划问题
\begin{gather}
	\max\left\{\sum_i Be^{-\beta t_i}\dfrac{C(t_i)^{1-\theta}}{1-\theta}\Delta t_i
	\right\}
\end{gather}
s.t.
\begin{gather}
	\sum_i e^{-R(t_i)}c(t_i)e^{(n+g)t_i}\Delta t_i\le
	k(0)+\sum_ie^{-R(t_i)}w(t_i)e^{(n+g)t_i}\Delta t_i
\end{gather}
从黎曼可积的角度来看，上述新规划问题等价于老规划问题。而新规划问题是典型的离散时间规划问题，因此可以采用拉格朗日方法进行求解。构造拉格朗日方程为
\begin{equation}
	\begin{aligned}
	& \max _{c\left(t_i\right)} \sum_i B e^{-\beta t_i} \frac{c\left(t_i\right)^{1-\theta}}{1-\theta} \Delta t_i+ \\
	& \lambda\left(-\sum_i e^{-R\left(t_i\right)} c\left(t_i\right) e^{(n+g) t_i} \Delta t_i+k(0)+\sum_i e^{-R\left(t_i\right)} w\left(t_i\right) e^{(n+g) t_i} \Delta t_i\right)
	\end{aligned}
	\end{equation}


其中$ \lambda $为预算约束方程所对应的拉格朗日乘子。拉格朗日函数关于消费的$ c(t_i) $一阶条件为
\begin{gather}
	Be^{^{-\beta t_i}}c(t_i)^{-\theta}\Delta t_i=\lambda e^{-R(t_i)}
	e^{(n+g)t_i}\Delta t_i
\end{gather}
消掉时间微分，得到
\begin{gather}
	Be^{-\beta t_i}c(t_i)^{-\theta}=\lambda e^{-R(t_i)}e^{(n+g)t_i}
\end{gather}
考虑到时间的任意性，得到
\begin{gather}
	Be^{-\beta t}c(t)^{-\theta}=\lambda e^{-R(t)}e^{(n+g)t}
\end{gather}
变形得到
\begin{gather}
	c(t)^{-\theta}=\dfrac{\lambda}{B}e^{\beta t-R(t)+(n+g)t}
\end{gather}
方程两边对时间求导，得到
\begin{gather}
	-\theta c(t)^{-\theta-1}\dot{c}(t)=\dfrac{\lambda}{B}e^{\beta t-R(t)+(n+g)t}(\beta-r(t)+n+g)
\end{gather}
上面两个方程相除，得到如下连续时间形式的欧拉方程（Euler equation）
\begin{align}
  \dfrac{\dot{c}(t)}{c(t)}=\dfrac{r(t)-\beta -n-g}{\theta}=\dfrac{r(t)-\rho-\theta g}{\theta}
\end{align}
第二个等号使用了定义$ \beta =\rho-n-(1-\theta) g $
还可以计算人均消费C的增长率
\begin{align}
\begin{split}
\dfrac{\dot {C}(t)}{C(t)}&=\dfrac{\dot{A}(t)}{A(t)}+\dfrac{\dot{c}(t)}{c(t)}
\\
&=g+\dfrac{r(t)-\rho-\theta g}{\theta}
\\
&=\dfrac{r(t)-\rho}{\theta}
\end{split}	
\end{align}
【注意：以上公式可以用来定义跨期替代弹性。】

还可以使用扰动法求得上述欧拉方程。家庭减少$ \Delta c $所损失的效用为$ Be^{-\beta t}c(t)^{-\theta}\Delta c $。家庭在$ t+\Delta t $时刻的消费可以增加$ e^{[r(t)-n-g]\Delta t}\Delta c $，并且
$ c(t+\Delta t) \approx c(t) e^{[\dot{c}(t)/c(t)]\Delta t}$。因此，可以得到

\[Be^{-\beta t}c(t)^{-\theta}\Delta c =Be^{-\beta (t+\Delta t)}\left(c(t)e^{[\dot{c}(t)/c(t)]\Delta t}\right)^{-\theta}e^{[r(t)-n-g]\Delta t} \Delta c\]


化简、变形得到
\begin{gather}
	\beta \Delta t+\theta\dfrac{\dot{c}(t)}{c(t)}\Delta=
	\left(r(t)-n-g
	\right)\Delta t
\end{gather}

消去时间微分并变形就得到上述欧拉方程。


\section{经济的动力系统} 
\subsection{c的动力系统} 
将$ r(t)=f'(k(t)) $代入到欧拉方程，得到消费的动力系统为

\begin{gather}
	\dfrac{\dot{c}(t)}{c(t)}=\dfrac{f'(k(t))-\rho-\theta g}{\theta}
\end{gather}
\begin{figure}[H]
	\centering
	\includegraphics[width=16cm,height=10cm]{2-1}
	\caption{FIGURE 2.1 The dynamics of c}
\end{figure}
注意：如何理解k较小时，消费增长率增加？\\
答：k较小，r就越大，从家庭欧拉方程的直观含义来看，资本的跨期期望收益就越高，因此消费增加就越高。


\subsection{k的动力系统} 
而资本的动态方程为
\begin{gather}
	\dot{k}(t)=f(k(t))-c(t)-(n+g)k(t)
\end{gather}
资本不变时，有$ c=f(k)-(n+g)k $, $ c(0)=f(0)-0=0 $, and $ d^2c/dk^2<0 $. 
\begin{figure}[H]
	\centering
	\includegraphics[width=15cm,height=12cm]{2-2}
	\caption{FIGURE 2.2 The dynamics of k}
\end{figure}
\subsection{相图} 
\begin{figure}[H]
	\centering
	\includegraphics[width=16cm,height=10cm]{2-3}
	\caption{FIGURE 2.3 The dynamics of c and k}
\end{figure}
问题：为什么有$ k* < k_{GR } $?\\
稳态资本为
\begin{gather}
	\dfrac{f'(k(t))-\rho-\theta g}{\theta}=0
\end{gather}
因此，
\begin{gather}
	f'(k*)=\rho+\theta g
\end{gather}
而黄金规则水平的资本由如下方式得到
$ max c=f(k)-(n+g)k $
得到
$ f'(k_{GR})=n+g $

我们已经假设了$ \rho-n-(1-\theta)g>0 $，所以$ \rho+\theta g>n+g $
而$ f''<0 $，因此，$ k* < k_{GR} $


\subsection{c的初始值} 
\begin{figure}[H]
	\centering
	\includegraphics[width=16cm,height=10cm]{2-4}
	\caption{The behavior of  c and k for various initial values of c}
\end{figure}

\begin{align}
\begin{cases}
\dfrac{\dot{c}(t)}{c(t)}=\dfrac{f'(k(t))-\rho-\theta g}{\theta}\\
\dot{k}(t)=f(k(t))-c(t)-(n+g)k(t)
\end{cases}	
\end{align}

假设k(0)<k*。
\begin{itemize}
\item A点：c(0)太高，$ \dot{c}>0,\dot{k}<0 $。
\item B点：c(0)处在$ \dot{k}=0 $上，所以$ \dot{c}>0 $，而$ \dot{k} $初始为0，经济体先向上移动，然后$ \dot{c} >0,\dot{k}<0$。
\item C点：$ c(0) $过高，所以虽然为正但是数值较小，而$ \dot{c}>0 $，因此，经济体初始向右上方向移动但$ k $右移幅度有限而无法到达$ E $点，当同$ \dot{k}=0 $轨迹相交后，$ \dot{c} >0,\dot{k}<0$。
\item D点：$ c(0) $过低。$ \dot{c},\dot{k} $初始均为正。由于$ \dot{c} $与$ c $成比例，当$ c $较小时，$ \dot{c} $也较小，$ c $太低因而上升幅度有限，当与$ \dot{c} =0$相交后，$ \dot{c} <0,\dot{k}>0$，经济向右下移动。
\item $ F $点：只有这一点才能到达均衡点$ E $
\end{itemize}




特别要注意的是，以上这些轨迹都满足两个一阶条件，但是我们还没有施加预算约束和资本不为负约束。
\begin{itemize}
\item 如果经济体开始于高于$ F $点，$ k $最终为负，因此可以排除这种情形。
\item 如果经济体开始于低于F点，如D点。最终$ k $会超过$ k_{GR} $，那么过了此时，利率$ r=f'(k)<n+g $，因此$ e^{-R(s)}e^{(n+g)s}=e^{\int_{0}^{s}[f^{\prime}(k)-(n+g) ]dt  } $随$ s $上升，因而$ e^{-R(s)}e^{(n+g)s}k(s) $ 发散。按照预算约束$ \lim\limits_{s\to \infty}e^{-R(s)} e^{(n+g)s}k(s)\ge 0$的含义，这说明：相对于家庭的终身消费的贴现值，其终身收入的贴现值是无穷大的，因此，在每一个时点，家庭可以增加消费从而获得更高的效用，这意味着家庭之前没有最大化其效用。因此，这不是均衡路径。
\item 如果经济体开始于$ F $点。$ k $收敛于$ k^* $，而$ r $收敛于$ f'(k^*) =\rho+\theta g$，因此，$ e^{-R(s)}e^{(n+g)s}=e^{-rs}e^{(n+g)s}=e^{-(\rho+\theta g)s+(n+g)s}=e^{-(\rho-n-(1-\theta )g)s}=e^{-(\rho- n-(1-\theta)g)s}=e^{-\beta s} $，其中$ \beta=\rho-n-(1-\theta)g $，$ e^{-R(s)} e^{(n+g)s}$以$ \beta $的速度下降，$ \lim\limits_{s\to \infty} e^{-R(s)}e^{(n+g)s}K(S)$为0。因此，由$ F $点开始的路径是唯一可行的路径。	
\end{itemize}


\subsection{鞍点路径}
鞍点路径：对于$ k $的任何初始值，存在一个唯一的$ c $的初始值。
\begin{figure}[H]
	\centering
	\includegraphics[width=16cm,height=10cm]{2-5}
	\caption{The saddle path}
\end{figure}
\section{福利}
 
\begin{itemize}
\item 福利经济学第一定理：如果市场是竞争性的、完全的，并且不存在外部性（以及如果行为者的人数是有限的），那么分散化的均衡是帕雷图最优的，也即不可能在不使某些人情境况变坏的前提下使某些人的境况变好。
\item 拉姆齐模型满足福利经济学第一定理的条件，因此均衡必定是帕雷图最优的。下面用中央计划者的规划问题来说明这一点。将会发现中央计划者的规划问题的最优解与分散的竞争性最优解，即拉姆齐模型的最优解，是相同的。	
\end{itemize}





中央计划者的目标是最大化家庭的福利水平，
\begin{gather}
	\max \left\{B\int_{t=0}^{\infty}e^{-\beta t}\dfrac{c(t)^{1-\theta}}{1-\theta}dt
	\right\}
\end{gather}
s.t.

\begin{gather}
	\dot{k}(t)=f(k(t))-c(t)-(n+g)k(t)
\end{gather}

将此连续时间规划问题离散化后，然后可以采用拉格朗日方法或者以资本作为状态变量的贝尔曼方程方法来进行求解。\\
方法1：拉格朗日方法\\
将中央计划者的规划问题离散化，得到\\
\begin{gather}
	\max\left\{\sum_{i}Be^{-\beta t_i}\dfrac{c(t_i)^{1-\theta}}{1-\theta}\Delta t_i
	\right\}
\end{gather}
预算约束方程为
\begin{gather}
	k(t_{i+1})-k(t_i)=[f(k(t_i))-c(t_i)-(n+g)k(t_i)
	]\Delta t_i
\end{gather}
构造如下拉格朗日方程，
\begin{gather}
	\max \sum_{i}Be^{-\beta t_i}\dfrac{c(t_i)^{1-\theta}}{1-\theta}\Delta t_i+
	\lambda_i\left\{-k(_{i+1})+k(t_i)+[f(k(t_i))-c(t_i)-(n+g)k(t_i)]\Delta t_i
	\right\}
\end{gather}

\begin{gather}
	\mathbb{L(c(t_i,\lambda_i))}\sum_{i}\left\{Be^{-\beta t_i}\dfrac{c(t_i)^{1-\theta}}{1-\theta}\Delta t_i+
	\lambda_i\left\{-k(_{i+1})+k(t_i)+[f(k(t_i))-c(t_i)-(n+g)k(t_i)]\Delta t_i
	\right\}\right\}
\end{gather}
此处$ \lambda_i $是预算约束方程所对应的拉格朗日乘子。

拉格朗日方程关于控制变量消费的一阶条件为
\begin{gather}
Be^{-\beta t_i}c(t_i)\Delta t_i=\lambda_i\Delta t_i
\end{gather}

关于资本的一阶条件为
\begin{gather}
  \lambda_i=\left\{1+\left[f'(k(t_{t+i}))-n-g
  \right]\Delta _{t+1}
  \right\}\lambda_{i+1}
\end{gather}
将两个一阶条件合并，得到
\begin{gather}
	Be^{-\beta t_i}c(t_i)^{-\theta}=\left\{
	1+\left[f'(k(t_{t+1}))-n-g
	\right]\Delta t_{i+1}
	\right\}Be^{-\beta t_{i+1}}c(t_{i+1})^{-\theta}
\end{gather}
变形得到
\begin{gather}
1=\left\{1+\left[f'(k(t_{i+1}))-n-g
\right]\Delta t_{i+1}
\right\}e^{-\beta \Delta t_{i+1}}\left(\dfrac{c(t_{i+1})}{c(t_i)}
\right)	^{-\theta}
\end{gather}
其中$ \Delta t_{i+1} = t_{i+1}-t_i$。变形，得到

\begin{gather}
  1=\left\{1+\left[f'(k(t_{i+1}))-n-g
  \right]\Delta t_{i+1}
  \right\}e^{-\beta \Delta t_{i+1}}\left(1+\dfrac{c(t_{i+1})-c(t_i)}{c(t_i)}
  \right)	^{-\theta}
\end{gather}
变形，得到
\begin{gather}
	1=\left\{1+\left[f'(k(t_{i+1}))-n-g
	\right]\Delta t_{i+1}
	\right\}e^{-\beta \Delta t_{i+1}}\left(1+\dfrac{c(t_{i+1})-c(t_i)}{\Delta t_{i+1}}\dfrac{\Delta(t_{i+1})}{c(t_i)}
	\right)	^{-\theta}
\end{gather}
那么，有
\begin{gather}
		1=\left\{1+\left[f'(k(t_{i+1}))-n-g
	\right]\Delta t_{i+1}
	\right\}e^{-\beta \Delta t_{i+1}}
	\left\{
	\left(
	1+\dfrac{c(t_{i+1})-c(t_i)}{\Delta t_{i+1}}
	\right)	\dfrac{c(t_i)}{c(t_{i+1})-c(t_i)}
	\right\}
	^{-\theta\left(\dfrac{c(t_{i+1})-c(t_i)}{\Delta t_{i+1}}\dfrac{\Delta _{t+1}}{c(t_i)}\right)}
\end{gather}
使用指数函数的性质，并考虑到时间的任意性，那么得到
\begin{gather}
	1=e^[f'(k(t))-n-g
	]\Delta t e^{-\beta t} e^{-\theta\dfrac{\dot{c}(t)}{c(t)}}\Delta t
\end{gather}
两边取对数，消去时间微分，得到
\begin{gather}
	0=f'(k(t))-n-g-\beta-\theta\dfrac{\dot{c}(t)}{c(t)}
\end{gather}
变形就得到了欧拉方程
\begin{gather}
	\dfrac{\dot{c}(t)}{c(t)}=\dfrac{f'(k(t))-\beta-n-g}{\theta}
\end{gather}

此欧拉方程与竞争性均衡中的欧拉方程是相同的。由于两个模型中的第二个方程，即预算约束方程，是相同的。因此，中央计划者的最优解和竞争性均衡的解是完全相同的。这就是拉姆齐增长模型中所展示的福利经济学第一定理。

方法2：贝尔曼方程方法
将中央计划者的规划问题离散化，并构造如下贝尔曼方程
\begin{gather}
	V(k(t_i))=\max
	\left\{
	u(c(t_i))+V(k(t_{i+1}))
	\right\}
\end{gather}

预算约束方程为

\begin{gather}
  k(t_{i+1})-k(t_i)=
  \left[
  f(k(t_i))-c(t_i)-(n+g)k(t_i)
  \right]\Delta t_i
\end{gather}

其中
\begin{gather}
	u(c(t_i))=Be^{-\beta t_i}\dfrac{c(t_i)^{1-\theta}}{1-\theta}\Delta t_i
\end{gather}

以上动态规划问题的控制变量是当期消费和下一期的资本存量，状态变量是当期的资本存量。此规划问题的特点是资本既是下一期的状态变量，也是当期的控制变量。使用预算约束方程将消费解出来，并代入效用函数，得到如下无约束的贝尔曼方程，
\begin{gather}
	V(k(t_i))=\max_{k(t_{i+1})}
	\left\{
	u\left(
	f(k(t_i))-(n+g)k(t_i)-\dfrac{k(t_{i+1})-k(t_i)}
	{\Delta t_i}
	\right)+V(k(t_{i+1}))
	\right\}
\end{gather}
显然，此规划问题的状态变量是当前资本，控制变量是下一期的资本。对控制变量求一阶条件，得到
\begin{gather}
	u'(c(t_i))\dfrac{1}{\Delta t_i}=V'(k(t_{i+1}))
\end{gather}
对状态变量使用包络定理，得到

\begin{gather}
    V'(k(t_{i}))=u'(c(t_i))
    \left(
    f'(k(t_{i+1}))-(n+g)+\dfrac{1}{\Delta t_i}
    \right)
\end{gather}
上述两个方程合并，得到

\begin{gather}
	u'(c(t_i))=\dfrac{1}{\Delta t_i}=
	u'(c(t_{i+1}))
	\left(
	f'(k(t_{i+1}))\Delta t_{i+1}-(n+g)\Delta t_{i+1}+1
	\right)
\end{gather}

而

\begin{gather}
	u'(c(t_i))\dfrac{1}{\Delta t_i}=
	B e^{-\beta t_i}c(t_i)^{-\theta}
\end{gather}

因此，

\begin{gather}
	e^{-\beta t_i} c(t_{i})^{-\theta}
	\left(
	1+f'(k(t_{i+1}))\Delta _{i+1}-(n+g)\Delta_{i+1}
	\right)
\end{gather}

我们在前面已经证明了上面的方程可以推出所要求取的欧拉方程。因此，使用贝尔曼方程方法和拉格朗日方法求得的结果是相同的。

\section{平衡增长路径}

\begin{itemize}
	\item 储蓄率为$ (y-c)/y $是内生的，但是不变的。
	\item 人均资本、人均产出、人均消费的增长率为$ g $。
	\item 即使储蓄率是内生的，劳动效率的增长仍然是人均产出持续增长的唯一源泉。
	
	\item 由于$ k^* < kGR $，所以被称为“修正的黄金规则资本存量”
\end{itemize}



\section{贴现率下降的效应} 
假设$ \rho $下降。变化与索洛模型中储蓄率上升十分相似。
\subsection{定性效应} 
\begin{figure}[H]
	\centering
	\includegraphics[width=16cm,height=10cm]{2-6}
	\caption{The effects of a fall in the discounts rate}
\end{figure}
\begin{align}
	\begin{cases}
	\dfrac{\dot{c}(t)}{c(t)}=\dfrac{f'(k(t))-\rho-\theta g}{\theta}\\
	\dot{k}(t)=f(k(t))-c(t)-(n+g)k(t)
	\end{cases}
\end{align}

\begin{itemize}
	\item 由第一个方程，如果$ \rho $降低，那么$ k^* $增加才能确保新均衡时$ dc/dt=0 $。
	\item $ \rho $变动时，$ k $不能跳跃，但是$ c $必须向下跳跃至新的鞍点路径才能收敛至平衡增长路径。
\end{itemize}

\subsection{调整速度和鞍点路径的斜率} 
对两个一阶条件
\begin{gather}
\begin{cases}
\dfrac{\dot(c)(t)}{c(t)}=\dfrac{f'(k(t))-\rho-\theta g}
{\theta}\\
\dot{k}(t)=f(k(t))-c(t)-(n+g)k(t)
\end{cases}
\end{gather}
，
在$ k=k^*$和$c=c^* $附近，取一阶泰勒近似，得到
\begin{gather}
	\dot{c}\simeq \dfrac{\partial \dot{c}}{\partial k}
	[k-k^*]+\dfrac{\partial \dot{c}}{\partial c}[c-c^*]
	\\
	\dot{k}\simeq \dfrac{\partial \dot{k}}{\partial k}
	[k-k^*]+\dfrac{\partial \dot{k}}{\partial c}[c-c^*]
\end{gather}

\begin{eqnarray}
	\dot{c-c^*}\sim \dfrac{f''(k^*)c^*}{\theta}(k-k^*)
	\\
	\dot{k-k^*}\approx [f''(k^*)-n-g](k-k^*)-(c-c^*)
	\\
	=[\rho+\theta g-n-g](k-k^*)-(c-c^*)
	\\
	=\beta(k-k^*)-(c-c^*)
\end{eqnarray}
变形，得到：
\begin{align}
\dfrac{\dot{c-c^*}}{c-c^*}\approx \dfrac{f''(k^*)c^*}{\theta}	\dfrac{k-k^*}{c-c^*}
\\
\dfrac{\dot{k-k^*}}{k-k^*}\approx\beta -\dfrac{c-c^*}{k-k^*}
\end{align}

因为是在$ k=k^* $和$ c=c^* $附近进行分析，所以可以考虑一种线性的情形：经济体沿着一条直线移向$ (k^*,c^*) $。在这种情形下：

\begin{itemize}
	\item 按照以上两个方程，消费和资本的运动速度是不变的。
	
	\item 经济体与$ (k^*,c^*) $的距离减少的速度不变，这样才能保证沿着直线移动。
	
\end{itemize}

总结一下：一阶近似取直线$ \Rightarrow $斜率$ \dfrac{k-k^*}{c-c^*} $是常数两者$ \Rightarrow $移动速度都是常数两者速度相等才能保证沿着直线移动
因此，设定：$ \dfrac{\dot{c-c^*}}{c-c^*}=\dfrac{\dot{k-k^*}}{k-k^*}\equiv  \mu$，那么，

\begin{gather}
	\mu=\dfrac{f''(k^*)c^*}{\theta}\dfrac{k-k^*}{c-c^*},and \, \mu=\beta-\dfrac{c-c^*}{k-k^*}\Rightarrow
	\\
	\mu=\beta-\dfrac{f''(k^*)c^*}{\theta}\dfrac{1}{\mu}
	\\
	\mu^2-\beta \mu+\dfrac{f''(k^*)c^*}{\theta}=0
\end{gather}

求解此方程，得到：
$ \mu=\dfrac{\beta\pm[\beta^2-4f''(k^*)c^*/\theta]^{1/2}}{2} $

$ \mu $为负才能确保距离减少，即走向$ (k^*,c^*) $，因此
\[ \mu_1=\{\beta-[\beta^2-4f''(k*)c^*/\theta]^{1/2}\}/2 \]

\begin{figure}[H]
	\centering
	\includegraphics[width=16cm,height=10cm]{2-7}
	\caption{The linearized phase diagram}
\end{figure}
\textbf{c和k的线性化动力系统：}\\
由于
\begin{gather}
	\mu_1=\dfrac{\dot{c-c}^*}{c-c^*}\approx \dfrac{f''(k^*)c^*}{\theta}\dfrac{k-k^*}{c-c^*}
\end{gather}
变形得到，
\begin{gather}
	c=c^*+\dfrac{f''(k^*)c^*}{\theta \mu_1}(k-k^*)
\end{gather}
因此，在$ 0 $时刻，c会跳跃至上述方程刻画的路径。然后，c和k均会以速度收敛至平衡增长路径，即
\begin{gather}
	k(t)=k^*+e^{\mu_1 t}[k(0)-k^*],c(t)=c^*+e^{\mu_1 t}[c(0)-c^*]
\end{gather}



\subsection{调整速度 }
由于
\begin{gather}
	\dfrac{\dot{c-c^*}}{c-c^*}=\dfrac{\dot{k-k^*}}{k-k^*}=\mu_1
\end{gather}

那么，
\begin{align}
\begin{split}
c(t)-c^*=e^{\mu_1 t}[c(0)-c^*],\\
k(t)-k^*=e^{\mu_1 t}[k(0)-k^*].
\end{split}
\end{align}

考虑C-D生产函数，$ f(k)=k^\alpha $。可以计算得到：
\begin{align}
	f''(k^*)=\alpha(\alpha-1)k^{*\alpha-2}=
	[(\alpha-1)/\alpha] r^{*2}/f(k^*)
\end{align}

此处$ r^*=\alpha k^{*\alpha-1} $是平衡增长路径上的实际利率。因此，

\begin{align}
	\mu_1=\dfrac{1}{2}\left\{\beta-
	\sqrt{
		\beta^2+\dfrac{4}{\theta}
		\dfrac{1-\alpha}{\alpha}
		r^{*2}(1-s^*)
	}
	\right\}
\end{align}

此处$ s^*=1-c^*/f(k^*) $是储蓄率。


在平衡增长路径上，储蓄为$ (n+g)k^* $（回忆预算约束方程均衡时有）$ \dot{k}(t) =f(k(t))-c(t)-(n+g)k(t)=0$$ \rightarrow $$ \texttt{储蓄}=f-c=(n+g)k(t) $, 因此
\begin{gather*}
s^*=(n+g)k^*/ k^{*\alpha}=\alpha(n+g)/r*	.
\end{gather*}

由$  \dfrac{\dot{c}(t)}{c(t)}=\dfrac{f'(k(t))-\rho-\theta g}{theta}$, 有$ r^*=\rho+\theta g $。将以上事实代入$ \mu_1 $ 表达式，得到

\begin{align}
	\mu_1=\dfrac{1}{2}\left\{
	\beta-
	\sqrt{
	\beta^2+\dfrac{4}{\theta}\dfrac{1-\alpha}{\alpha}
	(\rho+\theta g)(\rho+\theta g-\alpha(n+g))
}
	\right\}
\end{align}

参数校准：\\

取：$ \alpha=1/3,\rho=4\%,n=2\%,g=1\%,\theta=1. $\\
得到：$ r^*=5\%, s^*=20\%, \beta=2\%,\mu_1=-5.4\%. $

\begin{align}
	c(t)-c^*=e^{\mu_1 t}[c(0)-c^*],
	\\
	k(t)-k^*=e^{\mu_1 t}[k(0)-k^*].
\end{align}

以上结果与索洛模型相比,

\begin{align}
	\begin{split}
	k(t)-k^*=e^{-[1-\alpha_K](n+g+\delta)t}\left(k(0)-k^*\right)
	\\
	k(t)-k^*=e^{-2\% t[k(0)-k^*]}.
	\end{split}
\end{align}
【注意：此处模型没有折旧】

问题：之所以RCK模型收敛速度更快，是因为：当$ k<k^* $时，$ s>s* $；当$ k>k^* $时，$ s<s^* $，而索洛模型中使用了常数的储蓄率。

\textbf{解答：}

\begin{align}
	\dfrac{c_c^*}{k-k^*}=\beta-\mu_1
	\\
	\rightarrow c=c^*+(\beta-\mu_1)(k-k^*)
\end{align}

因此，储蓄率为

\begin{align}
	s=\dfrac{f(k)-c}{f(k)}
	\\
	\dfrac{f(k)-c^*-(\beta-\mu_1)(k-k^*)}{f(k)}
	\\
    =\dfrac{f(k^*)+f'(k^*)(k-k^*)-c^*-(\beta-\mu_1)(k-k^*)}{f(k)}
    \\
    \cdots >s^*,when \; k<k^*
\end{align}



\section{政府购买的效应}

\subsection{将政府纳入模型中}
假定政府在每单位时间以每单位有效劳动$ G(t) $的速度购买产出。假设政府购买不影响家庭效用。政府购买由总量税（lump-sum taxes）来融资，因此政府维持平衡预算。

资本的动力系统为：
\begin{gather}
	\dot{k}(t)=f(k(t))-c(t)-G(t)-(n+g)k(t)
\end{gather}

家庭最大化：

\begin{gather}
	B\int_{t=0}^{\infty}e^{-\beta t}\dfrac{c(t)^{1-\theta}}{1-\theta}dt
\end{gather}

s.t. 

\begin{gather}
	\int_{t=0}^{\infty}e^{-R(t)}c(t)e^{(n+g)t}dt\le 
	k(0)+\int_{t=0}^{\infty}e^{-R(t)}[w(t)-G(t)]e^{(n+g)t}dt
\end{gather}

得到的欧拉方程不变。而资本动力系统发生变化。
\begin{gather}
	\dfrac{\dot{c}(t)}{c(t)}=\dfrac{f'(k(t))-\rho-\theta g}{\theta}
	\\
	\dot{k}(t)=f(k(t))-c(t)-G(t)-(n+g)k(t)
\end{gather}

\subsection{政府购买的永久性和暂时性变动的效应} 

(a) 政府购买永久性增加

假设$ G(t) = G_L $，现在$ G $未预期的、永久性的增加到$ G_H $。由于不影响欧拉方程，所以轨迹不变。0的轨迹向下移动的数量等于G的增加数量。如下图。

\begin{figure}[H]
	\centering
	\includegraphics[width=16cm,height=10cm]{2-8}
	\caption{The effects of permanent increase in goverment purchases}\label{The effects of permanent increase in goverment purchases}
\end{figure}

\textbf{c如何跳跃？}\\
c下降的数量正好等于G增加的数量，并且资本存量与实际利率不变。

(b) 暂时性增加

\begin{itemize}
\item 假设G的增加是暂时性的：t0 时刻G增加，t1 时刻G恢复到GL。
\item c 不会下降GH-GL，家庭的效用函数为凹函数（concave utility function），更喜欢相对平滑的消费流。（回忆Jensen’s inequality） 
\item 由，而，所以dk/dt<0. 而r=f’(k)。因此，见下面的图（a）（b）。
\item 图（a）中G的增加相对持久。图（c）中G的增加相对短期。
\end{itemize}

\begin{figure}[H]
	\centering
	\includegraphics[width=16cm,height=12cm]{2-9}
	\caption{The effects of temporary inrease in goverment purchases}
\end{figure}

\subsection{经验应用：战争与实际利率} 

\begin{itemize}
	\item 前面表明：暂时性的政府购买增加导致实际利率上升。
	\item 暂时性政府购买增加的一个自然例子是战争。
\end{itemize}
	
\begin{figure}[H]
	\centering
	\includegraphics[width=15cm,height=10cm]{2-10}
	\caption{The temporary military spending and the long-term interest rate in the United Kindom (from Barro ,1987; used with permission)}
\end{figure}

Barro（1987）的估计结果：
\begin{figure}[H]
	\centering
	\includegraphics[width=15cm,height=4cm]{2-10(a)}
\end{figure}
其中$ R_t $ 长期名义利率（没有短期实际利率数据，而且样本期内预期通胀率没有发生系统性变动），是暂时性军事支出增加占GNP比重，
$ \widetilde{G}_{t} $是残差的一阶自回归模型的系数。
\[ R_t=a+bG_t+u_t, u_t=\rho_{t-1}+v_t \]
\begin{flalign}
	Cochrane－Orcutt:\\
	Step1: OLS: R_t=a+b\widetilde{G}_{t}+\widehat{\mu}_{t}\\
	Step2: ARMA:  \widehat{\mu}=\rho\widehat{\mu}_{t-1} +v_t and get \widehat{\rho}
	Step3: OLS:  R_t-\widehat{\rho}R_{t-1}=a(1-\widehat{\rho})+b(\widehat{G}_{t}-\widehat{\rho}\widehat{G}_{t-1})+v_t
	to get a and b.
\end{flalign}



Cochrane-Orcutt:\\
Step1: OLS: $R_t=a+b \tilde{G}_t+\hat{u}_t$\\
Step2: ARMA: $\hat{u}_t=\rho \hat{u}_{t-1}+v_t$ and get $\hat{\rho}$\\
Step3: OLS: $R_t-\hat{\rho} R_{t-1}=a(1-\hat{\rho})+b\left(\tilde{G}_t-\hat{\rho} \tilde{G}_{t-1}\right)+v_t$ to get $a$ and $b$.\\







































































